\documentclass[11pt]{article}
\usepackage[margin=1in]{geometry}
\usepackage{amsmath,amssymb}
\usepackage{booktabs}
\usepackage{enumitem}
\usepackage{xcolor}
\usepackage{titlesec}
\usepackage{hyperref}
\hypersetup{colorlinks=true,linkcolor=blue!50!black,citecolor=blue!50!black}

\titleformat{\section}{\normalfont\Large\bfseries\color{blue!35!black}}{\thesection}{0.6em}{}
\titleformat{\subsection}{\normalfont\large\bfseries\color{blue!25!black}}{\thesubsection}{0.6em}{}

\newcommand{\purpose}[1]{\par\noindent\textbf{Purpose:} #1\par}
\newcommand{\whenrun}[1]{\par\noindent\textbf{When it runs:} #1\par}
\newcommand{\decision}[1]{\par\noindent\textbf{Decision rule:} #1\par}

\title{\textbf{Statistical Testing \& Classification Reference}\\[4pt]
\large Updated Pipeline --- Detection and Attribution of AI-Generated and\\
Human-Written Phishing Emails}
\author{Mannan Bhardwaj (25BYB0013) \and Bhoumik Singh Rajput (25BYB0120)}
\date{September 2026}

\begin{document}
\maketitle

\section*{Scope of This Document}
This document explains every statistical test and classification model in the \emph{current} version
of the pipeline --- the version with the normality-testing branch removed (Kruskal--Wallis is used for
all continuous features, unconditionally) and the two-group Chi-Square test added for categorical
attribution features. It does not describe earlier design iterations. For each test: what it does, when
it runs, its formula, a fully worked numeric example, and how to read the result. The last section walks
one example email through the entire pipeline end to end.

\section{Quick Reference}

\begin{table}[h]
\centering
\renewcommand{\arraystretch}{1.3}
\begin{tabular}{p{2.7cm}p{4.3cm}p{3.3cm}p{2.6cm}}
\toprule
\textbf{Name} & \textbf{Purpose} & \textbf{Feature type} & \textbf{Runs on} \\
\midrule
Kruskal--Wallis $H$ & Does a feature differ across the 3 classes? & Continuous & Legit / Human / AI \\
Chi-Square (3-group) & Does a flag depend on class? & Categorical & Legit / Human / AI \\
Mann--Whitney $U$ & Does a feature separate the two phishing types? & Continuous & Human vs.\ AI only \\
Chi-Square (2-group) & Does a flag separate the two phishing types? & Categorical & Human vs.\ AI only \\
Benjamini--Hochberg FDR & Corrects all $p$-values for multiple testing & All & --- \\
Cohen's $d$ & How large is a continuous difference, really? & Continuous & --- \\
Cram\'er's $V$ & How strong is a categorical association, really? & Categorical & --- \\
\midrule
Model 1: Rank-Binned NB & Classifies using verified continuous features & Continuous & Every new email \\
Model 2: Categorical NB & Classifies using verified categorical features & Categorical & Every new email \\
Meta-Naive Bayes & Combines Model 1 \& Model 2 into one verdict & --- & Every new email \\
\bottomrule
\end{tabular}
\end{table}

The first seven rows run \textbf{once}, during training, to decide which features are trustworthy.
The last three rows run \textbf{every time} a new email needs to be classified, and never involve a
hypothesis test or a $p$-value.

\section{Phase 1: Feature-Selection Tests}
These four tests are run once on the training set. Each produces a $p$-value per feature; none of them
ever run again once the feature list is locked.

\subsection{Kruskal--Wallis $H$ Test}
\purpose{Tests whether a continuous feature's values differ across the three classes (Legitimate,
Human-Phishing, AI-Phishing), without assuming the data is normally distributed.}
\whenrun{On every continuous feature, comparing all three classes at once.}

\[
H = \frac{12}{N(N+1)} \sum_{j=1}^{k} \frac{R_j^2}{n_j} - 3(N+1)
\]
where $N$ is the total sample size, $k$ is the number of classes, $R_j$ is the sum of ranks in class
$j$, and $n_j$ is the size of class $j$.

\paragraph{Worked example --- feature: sentence perplexity variance (``burstiness'').}
Four emails sampled from each class:

\begin{center}
\begin{tabular}{lccccc}
\toprule
Legitimate & 2 & 3 & 4 & 5 & \\
Human-Phishing & 9 & 10 & 11 & 12 & \\
AI-Phishing & 5 & 6 & 7 & 8 & \\
\bottomrule
\end{tabular}
\end{center}

Pool all 12 values and rank them 1 (lowest) to 12 (highest). The value 5 appears twice (once in
Legitimate, once in AI-Phishing), so both instances share the average of ranks 4 and 5, i.e., rank 4.5.

\begin{center}
\begin{tabular}{lcccc}
\toprule
Class & Ranks & \multicolumn{2}{c}{Rank sum $R_j$} \\
\midrule
Legitimate & 1, 2, 3, 4.5 & \multicolumn{2}{c}{$R_1 = 10.5$} \\
Human-Phishing & 9, 10, 11, 12 & \multicolumn{2}{c}{$R_2 = 42$} \\
AI-Phishing & 4.5, 6, 7, 8 & \multicolumn{2}{c}{$R_3 = 25.5$} \\
\bottomrule
\end{tabular}
\end{center}

With $N=12$, $k=3$, $n_j=4$ for every class:
\[
H = \frac{12}{12 \times 13}\left(\frac{10.5^2}{4}+\frac{42^2}{4}+\frac{25.5^2}{4}\right) - 3(13)
= 0.0769 \times 631.125 - 39 \approx 9.55
\]
\decision{Compare $H$ to a chi-square distribution with $k-1=2$ degrees of freedom. The critical value
at $p=0.05$ is $5.99$. Since $9.55 > 5.99$, the difference is statistically significant --- burstiness
genuinely differs across the three classes, and the feature proceeds to the FDR/effect-size check.}

\subsection{Chi-Square Test (Three-Group)}
\purpose{Tests whether a categorical (yes/no) feature is associated with class membership.}
\whenrun{On every categorical/flag feature, comparing all three classes at once.}

\[
\chi^2 = \sum_{i=1}^{r}\sum_{j=1}^{c} \frac{(O_{ij}-E_{ij})^2}{E_{ij}}
\]
where $O_{ij}$ is the observed count in cell $(i,j)$ and $E_{ij} = \dfrac{\text{row total} \times
\text{column total}}{\text{grand total}}$ is the count expected if the feature were unrelated to class.

\paragraph{Worked example --- feature: contains an urgency keyword (Yes/No), 100 emails per class.}
\begin{center}
\begin{tabular}{lccc}
\toprule
 & Yes & No & Total \\
\midrule
Legitimate & 10 & 90 & 100 \\
Human-Phishing & 70 & 30 & 100 \\
AI-Phishing & 55 & 45 & 100 \\
\midrule
Total & 135 & 165 & 300 \\
\bottomrule
\end{tabular}
\end{center}
Expected count per cell (e.g., Legitimate/Yes) $=\frac{135\times100}{300}=45$; every class's expected
Yes/No split is $45/55$ since class is assumed irrelevant. Summing $(O-E)^2/E$ over all six cells gives
$\chi^2 \approx 78.8$.
\decision{With $\mathrm{df}=(3-1)(2-1)=2$, the critical value at $p=0.05$ is $5.99$. $78.8 \gg 5.99$:
urgency-keyword usage is strongly tied to class.}

\subsection{Mann--Whitney $U$ Test (Attribution)}
\purpose{Tests whether a continuous feature separates AI-Phishing from Human-Phishing specifically ---
a different, narrower question than the three-group tests above.}
\whenrun{On every continuous feature, restricted to the phishing-only subset (Human-Phishing vs.\
AI-Phishing).}

\[
U = n_1 n_2 + \frac{n_1(n_1+1)}{2} - R_1
\]

\paragraph{Worked example --- feature: urgency-keyword count per email.}
\begin{center}
\begin{tabular}{lccccc}
\toprule
Human-Phishing & 6 & 7 & 8 & 9 & 10 \\
AI-Phishing & 2 & 3 & 3 & 4 & 5 \\
\bottomrule
\end{tabular}
\end{center}
Pooling and ranking all 10 values (the tied pair of 3's each take the average rank 2.5): Human-Phishing
ranks are $6,7,8,9,10 \to R_1=40$; AI-Phishing ranks are $1,2.5,2.5,4,5 \to R_2=15$.
\[
U = (5)(5) + \frac{5(6)}{2} - 40 = 25+15-40 = 0
\]
\decision{$U$ ranges from $0$ (complete separation) to $n_1 n_2=25$ (complete overlap). $U=0$ here means
every Human-Phishing email outranked every AI-Phishing email on this feature --- the strongest possible
signal, and a very small $p$-value.}

\subsection{Chi-Square Test (Two-Group, Attribution)}
\purpose{The categorical counterpart to Mann--Whitney $U$: tests whether a categorical/flag feature
separates AI-Phishing from Human-Phishing specifically.}
\whenrun{On every categorical feature, restricted to the phishing-only subset.}

Same formula as Section 2.2, but with only two classes (rows) instead of three.

\paragraph{Worked example --- feature: contains a suspicious link (Yes/No), 100 emails per phishing type.}
\begin{center}
\begin{tabular}{lccc}
\toprule
 & Yes & No & Total \\
\midrule
Human-Phishing & 60 & 40 & 100 \\
AI-Phishing & 25 & 75 & 100 \\
\midrule
Total & 85 & 115 & 200 \\
\bottomrule
\end{tabular}
\end{center}
Expected Yes per group $=\frac{85\times100}{200}=42.5$; expected No $=57.5$. Summing $(O-E)^2/E$ over all
four cells gives $\chi^2 \approx 25.08$.
\decision{With $\mathrm{df}=1$, the critical value at $p=0.05$ is $3.84$. $25.08 \gg 3.84$: link
presence is strongly tied to \emph{which type} of phishing an email is, not just whether it's phishing
at all.}

\section{Multiple-Testing Correction and Effect Size}
Every test above only tells you whether a difference is \emph{probably real}. These two checks decide
whether it is \emph{real and worth keeping}.

\subsection{Benjamini--Hochberg FDR Correction}
\purpose{When many features are tested at once, some will look ``significant'' by pure chance. This
correction controls the expected proportion of false positives among everything ultimately kept.}
\whenrun{Once, across every raw $p$-value produced by every test above, jointly.}

\[
P_{(i)} \leq \frac{i}{m}\,Q \qquad (Q = 0.05)
\]
Sort all $m$ p-values ascending and rank them $i=1,\dots,m$. Find the largest $i$ whose $p$-value is
still below its own threshold $\frac{i}{m}Q$; keep every feature ranked at or below that $i$.

\paragraph{Worked example --- 10 features tested.}
\begin{center}
\begin{tabular}{cccc}
\toprule
Rank $i$ & $p$-value & Threshold $\frac{i}{10}(0.05)$ & Passes? \\
\midrule
1 & 0.001 & 0.005 & \checkmark \\
2 & 0.004 & 0.010 & \checkmark \\
3 & 0.006 & 0.015 & \checkmark \\
4 & 0.011 & 0.020 & \checkmark \\
5 & 0.019 & 0.025 & \checkmark \\
6 & 0.022 & 0.030 & \checkmark \\
7 & 0.030 & 0.035 & \checkmark \\
8 & 0.040 & 0.040 & \checkmark \\
9 & 0.060 & 0.045 & $\times$ \\
10 & 0.080 & 0.050 & $\times$ \\
\bottomrule
\end{tabular}
\end{center}
\decision{The largest passing rank is 8, so features ranked 1--8 are kept; ranks 9--10 are discarded,
even though a naive uncorrected 0.05 cutoff would have kept more of them.}

\subsection{Cohen's $d$}
\purpose{Measures how large a continuous feature's difference between two groups actually is, in
standardized units --- protects against large sample sizes making trivial differences look
``significant.''}
\whenrun{On every continuous feature that passed its p-value test, for every pairwise comparison
needed.}

\[
d = \frac{\bar{x}_1-\bar{x}_2}{s_{\text{pooled}}}, \qquad
s_{\text{pooled}} = \sqrt{\frac{(n_1-1)s_1^2+(n_2-1)s_2^2}{n_1+n_2-2}}
\]
Rule of thumb: $|d|\approx0.2$ small, $0.5$ medium, $0.8$ large. This pipeline requires $|d|\geq0.20$.

\paragraph{Worked example --- feature: average sentence length.}
Legitimate: $12,14,13,15,16$ (mean $14$); Phishing: $15,17,16,14,18$ (mean $16$). Both groups have
standard deviation $\approx1.58$, so pooled std $\approx1.58$.
\[
d = \frac{16-14}{1.58} \approx 1.27
\]
\decision{$1.27$ is a large effect --- Phishing emails use noticeably longer sentences than Legitimate
ones, not just detectably-but-trivially longer ones.}

\subsection{Cram\'er's $V$}
\purpose{The categorical counterpart to Cohen's $d$: measures how strong a categorical association
actually is, on a $0$-to-$1$ scale.}
\whenrun{On every categorical feature that passed its p-value test.}

\[
V = \sqrt{\frac{\chi^2}{n\times(\min(\text{rows,cols})-1)}}
\]
Rule of thumb (1 degree of freedom): $V\approx0.1$ weak, $0.3$ moderate, $0.5+$ strong.

\paragraph{Worked example --- reusing the suspicious-link feature from Section 2.4 ($\chi^2=25.08$,
$n=200$, $\mathrm{df}=1$).}
\[
V = \sqrt{\frac{25.08}{200\times1}} = \sqrt{0.1254} \approx 0.354
\]
\decision{A moderate-to-strong association --- link presence is a genuinely useful attribution signal,
not just a statistical technicality.}

\section{Phase 2: Classification Models}
\textbf{These are not hypothesis tests and do not produce $p$-values.} They run for every new email,
using only the parameters (bin edges, frequency tables, class priors) that Phase 1 already computed and
locked. Nothing here is recalculated per email.

\subsection{Model 1: Rank-Binned Naive Bayes}
\purpose{Classifies using continuous features that survived Kruskal--Wallis, by binning each feature's
value into intervals (e.g., Low/Medium/High) and looking up how often each class historically fell into
that bin.}

\[
P(x_i \in \text{Bin}_b \mid C_k) = \frac{N_{b,k}+\alpha}{N_k+\alpha B}, \qquad \alpha = 1 \text{ (Laplace smoothing)}
\]

\paragraph{Worked example --- feature: burstiness, using training counts (1{,}000 emails/class).}
\begin{center}
\begin{tabular}{lccc}
\toprule
 & Low & Medium & High \\
\midrule
Legitimate & 700 & 250 & 50 \\
Human-Phishing & 200 & 400 & 400 \\
AI-Phishing & 650 & 300 & 50 \\
\bottomrule
\end{tabular}
\end{center}
A new test email's burstiness falls in the \textbf{High} bin:
\[
P(\text{High}\mid\text{Legit})=\frac{51}{1003}\approx0.051,\quad
P(\text{High}\mid\text{Human})=\frac{401}{1003}\approx0.400,\quad
P(\text{High}\mid\text{AI})=\frac{51}{1003}\approx0.051
\]

\subsection{Model 2: Categorical Naive Bayes}
\purpose{Classifies using categorical features that survived Chi-Square, by looking up how often each
class historically showed that category value.}

\[
P(x_i=v\mid C_k) = \frac{\text{Count}(x_i=v \text{ in } C_k)+\alpha}{\text{Total in } C_k+\alpha K}
\]

\paragraph{Worked example --- feature: urgency keyword present (Yes/No), 1{,}000 emails/class.}
\begin{center}
\begin{tabular}{lcc}
\toprule
 & Yes & No \\
\midrule
Legitimate & 50 & 950 \\
Human-Phishing & 700 & 300 \\
AI-Phishing & 550 & 450 \\
\bottomrule
\end{tabular}
\end{center}
The same test email contains an urgency keyword (\textbf{Yes}):
\[
P(\text{Yes}\mid\text{Legit})\approx0.051,\quad P(\text{Yes}\mid\text{Human})\approx0.700,\quad
P(\text{Yes}\mid\text{AI})\approx0.550
\]

\subsection{Meta-Naive Bayes Probability Combiner}
\purpose{Combines Model 1's and Model 2's independent judgments into one final probability per class,
using Bayes' theorem --- not a simple average, so confident agreement between the two models is
amplified and disagreement is suppressed.}

\[
P(C_k \mid P_1,P_2) = \frac{P(C_k)\,P_1(C_k)\,P_2(C_k)}{\sum_j P(C_j)\,P_1(C_j)\,P_2(C_j)}
\]

With equal priors $P(C_k)=\tfrac{1}{3}$:
\[
\text{Legit: } \tfrac{1}{3}(0.051)(0.051)\approx0.00087 \quad
\text{Human: } \tfrac{1}{3}(0.400)(0.700)\approx0.0933 \quad
\text{AI: } \tfrac{1}{3}(0.051)(0.550)\approx0.00935
\]
Normalizing (dividing by the sum $\approx0.1035$):
\[
P(\text{Legit})\approx0.8\%, \qquad P(\text{Human})\approx90.2\%, \qquad P(\text{AI})\approx9.0\%
\]

\subsection{Final Two-Step Decision Rule}
\textbf{Step 1 (Legitimate vs.\ Phishing):}
\[
P(\text{Phishing}) = P(\text{Human})+P(\text{AI}) = 90.2\%+9.0\% = 99.2\% \;(\geq 0.50 \Rightarrow \text{Phishing})
\]
\textbf{Step 2 (Human vs.\ AI, only if Phishing):}
\[
P(\text{AI}\mid\text{Phishing}) = \frac{9.0}{9.0+90.2} \approx 9.1\% \;(< 0.50 \Rightarrow \text{Human-Written})
\]

\section{End-to-End Summary}
One email, carried through every stage of the current pipeline:

\begin{enumerate}[leftmargin=1.4cm]
\item \textbf{Training time (once):} Kruskal--Wallis confirmed burstiness differs across classes
($H\approx9.55$); Chi-Square confirmed urgency-keyword use differs across classes ($\chi^2\approx78.8$).
Both passed FDR correction and their effect-size checks, so both were locked into the selected feature
set along with their trained bin edges / frequency tables.
\item \textbf{New email arrives:} its burstiness value lands in the \emph{High} bin; it contains an
urgency keyword.
\item \textbf{Classification:} Model 1 (Rank-Binned NB) and Model 2 (Categorical NB) independently score
the email against the stored training statistics; Meta-Naive Bayes combines both scores.
\item \textbf{Decision:} $P(\text{Phishing})=99.2\%\Rightarrow$ Phishing; $P(\text{AI}\mid\text{Phishing})
\approx9.1\%\Rightarrow$ \textbf{Human-Written}.
\end{enumerate}

Every number in this final verdict traces back to a specific, auditable test result from Phase 1 and a
specific, auditable probability lookup from Phase 2 --- which is exactly the ``full auditability''
property the pipeline is designed around.

\end{document}
